% Generated by Listas/make_statements.py from pratica-lecture-03.md, item 8. Do not edit.
\textbf{Source.} Tadelis (2013), \emph{Game Theory: An Introduction}, Exercise 7.4 ``Centipedes'', p.~148 {[}PDF 167{]}. Reproduced verbatim. Answer in the source: none in the book (the preface, p.~xiii, says about half are solved on Princeton UP's website; not checked).



\begin{quote}

\textbf{Centipedes:} Imagine a two-player game that proceeds as follows. A pot of money is created with \$6 in it initially. Player 1 moves first, then player 2, then player 1 again, and finally player 2 again. At each player's turn to move, he has two possible actions: grab (\(G\)) or share (\(S\)). If he grabs he gets \(\frac{2}{3}\) of the current pot of money, the other player gets \(\frac{1}{3}\) of the pot, and the game ends. If he shares then the size of the current pot is multiplied by \(\frac{3}{2}\) and the next player gets to move. At the last stage at which player 2 moves, if he chooses \(S\) then the pot is still multiplied by \(\frac{3}{2}\), player 2 gets \(\frac{1}{3}\) of the pot, and player 1 gets \(\frac{2}{3}\) of the pot.



\begin{enumerate}

\def\labelenumi{\alph{enumi}.}

\item

  Model this as an extensive-form game tree. Is it a game of perfect or imperfect information?

\item

  How many terminal nodes does the game have? How many information sets?

\item

  How many pure strategies does each player have?

\item

  Find the Nash equilibria of this game. How many outcomes can be supported in equilibrium?

\item

  Now imagine that at the last stage at which player 2 moves, if he chooses to share then the pot is equally split among the players. Does your answer to part (d) change?

\end{enumerate}

\end{quote}

